% Options for packages loaded elsewhere
\PassOptionsToPackage{unicode}{hyperref}
\PassOptionsToPackage{hyphens}{url}
\documentclass[
  english,
]{article}
\usepackage{xcolor}
\usepackage{amsmath,amssymb}
\setcounter{secnumdepth}{5}
\usepackage{iftex}
\ifPDFTeX
  \usepackage[T1]{fontenc}
  \usepackage[utf8]{inputenc}
  \usepackage{textcomp} % provide euro and other symbols
\else % if luatex or xetex
  \usepackage{unicode-math} % this also loads fontspec
  \defaultfontfeatures{Scale=MatchLowercase}
  \defaultfontfeatures[\rmfamily]{Ligatures=TeX,Scale=1}
\fi
% \usepackage{lmodern}
\ifPDFTeX\else
  % xetex/luatex font selection
\fi
% Use upquote if available, for straight quotes in verbatim environments
\IfFileExists{upquote.sty}{\usepackage{upquote}}{}
\IfFileExists{microtype.sty}{% use microtype if available
  \usepackage[]{microtype}
  \UseMicrotypeSet[protrusion]{basicmath} % disable protrusion for tt fonts
}{}
\makeatletter
\@ifundefined{KOMAClassName}{% if non-KOMA class
  \IfFileExists{parskip.sty}{%
    \usepackage{parskip}
  }{% else
    \setlength{\parindent}{0pt}
    \setlength{\parskip}{6pt plus 2pt minus 1pt}}
}{% if KOMA class
  \KOMAoptions{parskip=half}}
\makeatother
\usepackage{color}
\usepackage{fancyvrb}
\newcommand{\VerbBar}{|}
\newcommand{\VERB}{\Verb[commandchars=\\\{\}]}
\DefineVerbatimEnvironment{Highlighting}{Verbatim}{commandchars=\\\{\}}
% Add ',fontsize=\small' for more characters per line
\newenvironment{Shaded}{}{}
\newcommand{\AlertTok}[1]{\textcolor[rgb]{1.00,0.00,0.00}{\textbf{#1}}}
\newcommand{\AnnotationTok}[1]{\textcolor[rgb]{0.38,0.63,0.69}{\textbf{\textit{#1}}}}
\newcommand{\AttributeTok}[1]{\textcolor[rgb]{0.49,0.56,0.16}{#1}}
\newcommand{\BaseNTok}[1]{\textcolor[rgb]{0.25,0.63,0.44}{#1}}
\newcommand{\BuiltInTok}[1]{\textcolor[rgb]{0.00,0.50,0.00}{#1}}
\newcommand{\CharTok}[1]{\textcolor[rgb]{0.25,0.44,0.63}{#1}}
\newcommand{\CommentTok}[1]{\textcolor[rgb]{0.38,0.63,0.69}{\textit{#1}}}
\newcommand{\CommentVarTok}[1]{\textcolor[rgb]{0.38,0.63,0.69}{\textbf{\textit{#1}}}}
\newcommand{\ConstantTok}[1]{\textcolor[rgb]{0.53,0.00,0.00}{#1}}
\newcommand{\ControlFlowTok}[1]{\textcolor[rgb]{0.00,0.44,0.13}{\textbf{#1}}}
\newcommand{\DataTypeTok}[1]{\textcolor[rgb]{0.56,0.13,0.00}{#1}}
\newcommand{\DecValTok}[1]{\textcolor[rgb]{0.25,0.63,0.44}{#1}}
\newcommand{\DocumentationTok}[1]{\textcolor[rgb]{0.73,0.13,0.13}{\textit{#1}}}
\newcommand{\ErrorTok}[1]{\textcolor[rgb]{1.00,0.00,0.00}{\textbf{#1}}}
\newcommand{\ExtensionTok}[1]{#1}
\newcommand{\FloatTok}[1]{\textcolor[rgb]{0.25,0.63,0.44}{#1}}
\newcommand{\FunctionTok}[1]{\textcolor[rgb]{0.02,0.16,0.49}{#1}}
\newcommand{\ImportTok}[1]{\textcolor[rgb]{0.00,0.50,0.00}{\textbf{#1}}}
\newcommand{\InformationTok}[1]{\textcolor[rgb]{0.38,0.63,0.69}{\textbf{\textit{#1}}}}
\newcommand{\KeywordTok}[1]{\textcolor[rgb]{0.00,0.44,0.13}{\textbf{#1}}}
\newcommand{\NormalTok}[1]{#1}
\newcommand{\OperatorTok}[1]{\textcolor[rgb]{0.40,0.40,0.40}{#1}}
\newcommand{\OtherTok}[1]{\textcolor[rgb]{0.00,0.44,0.13}{#1}}
\newcommand{\PreprocessorTok}[1]{\textcolor[rgb]{0.74,0.48,0.00}{#1}}
\newcommand{\RegionMarkerTok}[1]{#1}
\newcommand{\SpecialCharTok}[1]{\textcolor[rgb]{0.25,0.44,0.63}{#1}}
\newcommand{\SpecialStringTok}[1]{\textcolor[rgb]{0.73,0.40,0.53}{#1}}
\newcommand{\StringTok}[1]{\textcolor[rgb]{0.25,0.44,0.63}{#1}}
\newcommand{\VariableTok}[1]{\textcolor[rgb]{0.10,0.09,0.49}{#1}}
\newcommand{\VerbatimStringTok}[1]{\textcolor[rgb]{0.25,0.44,0.63}{#1}}
\newcommand{\WarningTok}[1]{\textcolor[rgb]{0.38,0.63,0.69}{\textbf{\textit{#1}}}}
\usepackage{longtable,booktabs,array}
\newcounter{none} % for unnumbered tables
\usepackage{calc} % for calculating minipage widths
% Correct order of tables after \paragraph or \subparagraph
\usepackage{etoolbox}
\makeatletter
\patchcmd\longtable{\par}{\if@noskipsec\mbox{}\fi\par}{}{}
\makeatother
% Allow footnotes in longtable head/foot
\IfFileExists{footnotehyper.sty}{\usepackage{footnotehyper}}{\usepackage{footnote}}
\makesavenoteenv{longtable}
\ifLuaTeX
\usepackage[bidi=basic,shorthands=off]{babel}
\else
\usepackage[bidi=default,shorthands=off]{babel}
\fi
\ifLuaTeX
  \usepackage{selnolig} % disable illegal ligatures
\fi
\setlength{\emergencystretch}{3em} % prevent overfull lines
\providecommand{\tightlist}{%
  \setlength{\itemsep}{0pt}\setlength{\parskip}{0pt}}
\usepackage[]{natbib}
\bibliographystyle{plainnat}
\makeatletter
\@ifpackageloaded{subcaption}{}{\usepackage{subcaption}}
\@ifpackageloaded{caption}{}{\usepackage{caption}}
\captionsetup[subfigure]{margin=0.5em}
\AtBeginDocument{%
\renewcommand*\figurename{Figure}
\renewcommand*\tablename{Table}
}
\AtBeginDocument{%
\renewcommand*\listfigurename{List of Figures}
\renewcommand*\listtablename{List of Tables}
}
\newcounter{pandoccrossref@subfigures@footnote@counter}
\newenvironment{pandoccrossrefsubfigures}{%
\setcounter{pandoccrossref@subfigures@footnote@counter}{0}
\begin{figure}\centering%
\gdef\global@pandoccrossref@subfigures@footnotes{}%
\DeclareRobustCommand{\footnote}[1]{\footnotemark%
\stepcounter{pandoccrossref@subfigures@footnote@counter}%
\ifx\global@pandoccrossref@subfigures@footnotes\empty%
\gdef\global@pandoccrossref@subfigures@footnotes{{##1}}%
\else%
\g@addto@macro\global@pandoccrossref@subfigures@footnotes{, {##1}}%
\fi}}%
{\end{figure}%
\addtocounter{footnote}{-\value{pandoccrossref@subfigures@footnote@counter}}
\@for\f:=\global@pandoccrossref@subfigures@footnotes\do{\stepcounter{footnote}\footnotetext{\f}}%
\gdef\global@pandoccrossref@subfigures@footnotes{}}
\@ifpackageloaded{float}{}{\usepackage{float}}
\floatstyle{ruled}
\@ifundefined{c@chapter}{\newfloat{codelisting}{h}{lop}}{\newfloat{codelisting}{h}{lop}[chapter]}
\floatname{codelisting}{Listing}
\newcommand*\listoflistings{\listof{codelisting}{List of Listings}}
\makeatother
\usepackage{bookmark}
\IfFileExists{xurl.sty}{\usepackage{xurl}}{} % add URL line breaks if available
\urlstyle{same}
% fallback for those not using the hyperref driver hyperxmp:
\makeatletter
\@ifundefined{xmpquote}{\newcommand{\xmpquote}[1]{#1}}{}
\makeatother
\hypersetup{
  pdftitle={template\_pitch\_deck: Reproducible{,} Validated Pitch-Deck Generation},
  pdfauthor={Daniel Ari Friedman},
  pdflang={en},
  colorlinks=true,linkcolor=red,urlcolor=red,citecolor=red,anchorcolor=red,filecolor=red,
  pdfcreator={LaTeX via pandoc}}

\title{template\_pitch\_deck: Reproducible, Validated Pitch-Deck
Generation}
\author{Daniel Ari Friedman}
\date{}

\IfFileExists{seqsplit.sty}{\usepackage{seqsplit}}{\newcommand{\seqsplit}[1]{#1}}
\protected\def\breakseq#1{\seqsplit{#1}}
\protected\def\breaktt#1{\begingroup\ttfamily\seqsplit{#1}\endgroup}
\title{template\_pitch\_deck: Reproducible, Validated Pitch-Deck Generation}
\author{Daniel Ari Friedman\\\footnotesize{Active Inference Institute}\\\footnotesize{\texttt{daniel@activeinference.institute}}\\\footnotesize{\href{https://orcid.org/0000-0001-6232-9096}{ORCID: 0000-0001-6232-9096}} \\ \footnotesize{DOI: \href{https://doi.org/10.5281/zenodo.21281509}{10.5281/zenodo.21281509}}}
\date{2026-07-08}

% Core mathematics
\usepackage{amsmath}
\usepackage{amssymb}

% Document layout
\usepackage{geometry}
\geometry{margin=1in}
\usepackage{float}
\usepackage{graphicx}

% Tables
\usepackage{booktabs}
\usepackage{longtable}

% Typography and formatting
\usepackage{microtype}
\usepackage{xcolor}

% Cross-references and citations
\usepackage{hyperref}
\hypersetup{
    colorlinks=true,
    allcolors=red
}
\usepackage[capitalise,noabbrev]{cleveref}
\usepackage{natbib}

\begin{document}
\begin{titlepage}
\centering
\vspace*{2cm}
{\LARGE\bfseries template\_pitch\_deck: Reproducible, Validated Pitch-Deck Generation\par}
\vspace{0.75em}
{\large Short/Medium/Long PDF+PPTX decks from one token-resolved content source\par}
\vfill
\makeatletter
{\@author\par}
\vspace{1em}
{\@date\par}
\makeatother
\vfill
\end{titlepage}
\thispagestyle{empty}
\tableofcontents
\newpage

\section{Abstract}\label{sec:abstract}

Research groups routinely need to pitch their work --- to funders,
partners, or collaborators --- yet pitch decks are almost never treated
as reproducible research artifacts: they are hand-assembled in
proprietary slide tools, contain unverifiable claims, and cannot be
regenerated when the underlying facts change.
\breaktt{template\_pitch\_deck} closes that gap. It generates six
artifacts from one token-resolved content source --- short, medium, and
long decks, each in both PDF and PPTX --- with every numeric claim
traced back to a live introspection of the repository it describes,
every \texttt{\{\{TOKEN\}\}} substitution verified to have actually
landed, and every sentence checked against a denylist of pitch-deck
clichés.

The flagship content pitches
\href{../../template_template/}{\breaktt{template\_template}}, this
monorepo's own autopoietic meta-project, to a meta-science and
science-integrity audience --- the kind of pitch an organization like
the Active Inference Institute or COGSEC would actually hand to a
funder. The rendering engine itself is new, reusable infrastructure:
\breaktt{infrastructure/rendering/slide\_deck.py} (ReportLab, PDF) and
\breaktt{infrastructure/rendering/pptx\_deck.py} (python-pptx) both
consume the same \texttt{DeckContent} model, so a PDF and a PPTX built
from identical content carry identical slide counts and identical text
--- verified by direct read-back of both file formats, not by
inspection.

\textbf{Keywords:} pitch deck, slide generation, reproducible research
communication, meta-science infrastructure, token validation, PPTX, PDF
rendering.

\newpage

\section{Introduction}\label{sec:introduction}

Every other public exemplar in this monorepo renders a manuscript:
prose, figures, citations, a combined PDF. None of them render a
\emph{pitch} --- a short-form, persuasion-oriented artifact whose job is
not to document a method but to move an audience to a decision. Pitch
decks live outside the repository's reproducibility guarantees entirely:
built by hand in proprietary tools, their claims untethered from any
generator, their content unchecked for the boilerplate language
(``synergy,'' ``disruptive,'' ``10x'') that makes a real pitch land
badly with a sophisticated audience.

\breaktt{template\_pitch\_deck} treats a pitch deck the same way this
repository treats a manuscript: as a build artifact with a single source
of truth, a validation gate, and a reproducibility guarantee. Three
properties make that possible.

\textbf{One content source, six artifacts.}
\texttt{manuscript/deck\_content\_\{short,medium,long\}.yaml} define the
slide-by-slide narrative at three lengths; the live token builder
(\breaktt{src/deck\_tokens.py}) supplies the facts.
\breaktt{scripts/render\_decks.py} resolves tokens once and calls both
renderers --- \breaktt{infrastructure.rendering.slide\_deck.render\_pdf}
and \breaktt{infrastructure.rendering.pptx\_deck.render\_pptx} ---
against the identical resolved content, so PDF and PPTX cannot drift
from each other.

\textbf{Facts, not fabrication.} Every numeric claim produced by
\breaktt{src/deck\_tokens.py} about the pitch subject
(\breaktt{template\_template}) is generated from a live read of that
project's own \texttt{README.md}/\texttt{AGENTS.md} or the repository's
public exemplar roster --- never hand-typed, never a plausible-sounding
invented metric. \breaktt{src/token\_resolution.py} raises loudly if any
\texttt{\{\{TOKEN\}\}} in the content source has no corresponding
resolved value, and a rendered artifact's extracted text is checked to
contain zero leftover \texttt{\{\{} literals.

\textbf{Cliché is a lint, not a vibe.} \breaktt{src/cliche\_lint.py} runs
a word-boundary denylist of pitch-deck stock phrases over every resolved
slide before render. A pitch that reads as generic --- regardless of how
factually accurate it is --- fails the same way an uncovered line of
code fails a coverage gate.

The remainder of this manuscript covers the deck-rendering architecture
(sec.~\ref{sec:architecture}), the validation model
(sec.~\ref{sec:validation}), and the reproducibility guarantees
(sec.~\ref{sec:reproducibility}) that let
\breaktt{uv\ run\ python\ scripts/render\_decks.py\ -\/-project\ templates/template\_pitch\_deck}
produce the same six files, byte-for-byte, on any machine with this repo
checked out.

\newpage

\section{Deck Rendering Architecture}\label{sec:architecture}

\subsection{Content model}\label{content-model}

\breaktt{infrastructure/rendering/slide\_deck.py} defines the
format-agnostic content model shared by both renderers:

\begin{itemize}
\tightlist
\item
  \texttt{Slide} --- one slide's title, bullets, an optional
  speaker-notes string, an optional local figure path, and a
  \texttt{kind} (\texttt{title}, \texttt{section}, or \texttt{content}).
\item
  \texttt{DeckContent} --- a deck title/subtitle plus an ordered tuple
  of \texttt{Slide}.
\item
  \texttt{SlideBudget} --- the three published lengths (\texttt{SHORT} \texorpdfstring{\ensuremath{\leq}}{<=}
  10 slides, \texttt{MEDIUM} \texorpdfstring{\ensuremath{\leq}}{<=} 22, \texttt{LONG} \texorpdfstring{\ensuremath{\leq}}{<=} 45) and
  \breaktt{filter\_deck\_for\_budget}, a pure function that truncates a
  deck to a budget without mutating slide order or content.
\end{itemize}

\subsection{Two renderers, one model}\label{two-renderers-one-model}

\breaktt{infrastructure/rendering/slide\_deck.render\_pdf} draws each
slide directly with ReportLab's canvas API onto a 16:9 page --- no
LaTeX/pandoc dependency, following the same ``project needs pixel-level
layout control'' reasoning that led \breaktt{template\_newspaper} to a
hand-written ReportLab engine, but implemented as a generic,
project-agnostic \emph{format} renderer (like
\breaktt{infrastructure/rendering/docx\_renderer.py} or
\breaktt{epub\_renderer.py}) rather than a project-owned layout engine,
because a slide deck --- unlike a newspaper's column geometry --- is a
reusable output shape any future project can consume.

\breaktt{infrastructure/rendering/pptx\_deck.render\_pptx} builds the
identical slide sequence with \texttt{python-pptx} (an opt-in
dependency: \breaktt{uv\ sync\ -\/-group\ rendering-pptx}), matching
title/section/content slide handling 1:1 with the PDF path. Both
renderers are exercised by
\breaktt{tests/infra\_tests/rendering/test\_slide\_deck.py} and
\breaktt{test\_pptx\_deck.py}, including a direct parity test that
renders the same \texttt{DeckContent} through both paths and asserts the
PDF page count equals the PPTX slide count.

\subsection{Project-side content}\label{project-side-content}

\breaktt{projects/templates/template\_pitch\_deck/src/} holds only
pitch-deck-domain code --- loading
\breaktt{manuscript/deck\_content\_*.yaml} into \texttt{DeckContent}
objects, resolving \texttt{\{\{TOKEN\}\}} values, and linting for cliché
--- never layout/drawing code. \breaktt{scripts/render\_decks.py} is the
thin orchestrator that ties content loading, token resolution, cliché
linting, and the two infra renderers together into the six published
artifacts under \texttt{output/\{pdf,pptx\}/}.

\newpage

\section{Content and Validation}\label{sec:validation}

\subsection{Methods}\label{methods}

\subsection{The flagship pitch}\label{the-flagship-pitch}

The shipped content pitches
\href{../../template_template/}{\breaktt{template\_template}} --- this
monorepo's autopoietic meta-project, which introspects the repository's
own architecture and regenerates its manuscript from live counts --- to
a meta-science and science-integrity audience. Every length follows the
same arc: the problem (research communication and reproducibility are
under-tooled), the solution (a two-layer, thin-orchestrator monorepo
with a 90\%/60\% coverage floor, no-mocks testing, and multi-platform
publication), proof (real, currently-measured facts: exemplar count,
coverage floors, the publishing surface \breaktt{template\_template}
itself already reaches), and an ask. Medium and long variants add
landscape, architecture, and roadmap detail; long adds a full
governance/confidentiality walkthrough and an appendix.

\subsection{Token resolution}\label{token-resolution}

\breaktt{src/token\_resolution.py} implements the same
\texttt{\{\{TOKEN\}\}} convention used by
\breaktt{infrastructure/rendering/manuscript\_injection.py}
(\texttt{\textbackslash{}\{\textbackslash{}\{{[}A-Z{]}{[}A-Z0-9\_{]}*\textbackslash{}\}\textbackslash{}\}}),
scoped to \breaktt{manuscript/deck\_content\_*.yaml} instead of
\texttt{manuscript/*.md}. \texttt{resolve\_tokens} raises if any token
in the content has no matching key in the live token set from
\breaktt{src/deck\_tokens.py} --- mirroring \texttt{template\_madlib}'s
\breaktt{test\_all\_manuscript\_tokens\_are\_generated} pre-substitution
coverage check, adapted to deck content.
\breaktt{scripts/audit\_deck\_content.py} runs this check plus the cliché
lint in one pass and exits non-zero on either class of failure; both
failure modes are proven to actually fire (a deliberately-broken fixture
triggers each) before the real content is checked clean.

\subsection{Cliché lint}\label{clichuxe9-lint}

\breaktt{src/cliche\_lint.py} maintains a word-boundary-safe denylist of
pitch-deck stock phrases (``synergy,'' ``disrupt,'' ``10x,''
``game-changing,'' ``rocket ship,'' ``paradigm shift,'' and more). It is
checked against every resolved slide across all three lengths as part of
the same audit script, and --- like the token check --- is proven to
fire on a deliberately cliché-laden sentence before the real content's
cleanliness is trusted.

\newpage

\section{Reproducibility}\label{sec:reproducibility}

\breaktt{scripts/render\_decks.py} is deterministic: given the same
repository state (\breaktt{manuscript/deck\_content\_*.yaml}, the live
tokens from \breaktt{src/deck\_tokens.py}, and the live facts
\breaktt{src/deck\_tokens.py} reads from \breaktt{template\_template}'s
own files and the public exemplar roster), two consecutive runs produce
byte-identical PDF and PPTX output. No wall-clock timestamps appear in
slide content; the only place a generation time is recorded is deck
metadata, following the \breaktt{STEGANOGRAPHY\_DETERMINISTIC}-style
convention already used elsewhere in this repository for reproducible
builds.

\begin{Shaded}
\begin{Highlighting}[]
\ExtensionTok{uv}\NormalTok{ run python scripts/render\_decks.py}
\CommentTok{\# → output/pdf/template\_template\_pitch\_\{short,medium,long\}.pdf}
\CommentTok{\# → output/pptx/template\_template\_pitch\_\{short,medium,long\}.pptx}

\ExtensionTok{uv}\NormalTok{ run pytest projects/templates/template\_pitch\_deck/tests/ }\DataTypeTok{\textbackslash{}}
  \AttributeTok{{-}{-}cov}\OperatorTok{=}\NormalTok{projects/templates/template\_pitch\_deck/src }\AttributeTok{{-}{-}cov{-}fail{-}under}\OperatorTok{=}\NormalTok{90}

\ExtensionTok{uv}\NormalTok{ run python scripts/audit/check\_template\_drift.py}
\end{Highlighting}
\end{Shaded}

This project participates in the standard multi-project pipeline like
any other public exemplar
(\breaktt{./run.sh\ -\/-project\ templates/template\_pitch\_deck\ -\/-pipeline\ -\/-core-only})
and requires no LLM/Ollama stage --- deck content is authored and
token-resolved, not generated at render time by a model call, which
keeps the artifact deterministic and the core pipeline Ollama-optional.

\newpage



\clearpage

\bibliography{references}
\end{document}
